\documentclass[conference]{IEEEtran}
\usepackage{amsmath,amssymb}
\usepackage{booktabs}
\usepackage{graphicx}
\usepackage{xcolor}
\usepackage{listings}
\usepackage{url}
\newcommand{\TBD}[1]{\textcolor{red}{[TBD: #1]}}
\newcommand{\sys}{\textsc{OpenSNAP}}
\newcommand{\ovql}{\textsc{OVQL}}
\lstset{basicstyle=\ttfamily\footnotesize,columns=fullflexible,keepspaces=true}

\begin{document}
\title{OpenSNAP: Open-Vocabulary Querying over 3D Point Clouds}
\author{\IEEEauthorblockN{\TBD{authors}}\IEEEauthorblockA{\TBD{affiliations}}}
\maketitle

\begin{abstract}
3D point clouds are produced at scale by indoor scanners and LiDAR, yet they are stored and exchanged as raw
coordinates: answering even simple questions such as \emph{``where are the fire extinguishers?''} or
\emph{``which chairs are next to a table?''} requires a bespoke perception pipeline per question and per vocabulary.
We present \sys{}, a system that materialises a point cloud into an \emph{object relation} with open-vocabulary,
probabilistic semantics and answers declarative queries over it. Objects are tuples carrying geometry, a quality
estimate and two language-aligned embeddings; no label column is stored, and the probability that an object matches
an arbitrary phrase is computed at query time and calibrated against a background vocabulary. \ovql{}, a small query
language, combines open-vocabulary predicates, pointing (click) predicates and spatial predicates under
tuple-independence semantics. When a query mentions something no stored object explains, \sys{} \emph{refines the
relation at query time}: it uses the query text as the click oracle of a promptable segmenter, materialises the missing
object and persists it --- database cracking applied to semantics. On ScanNet, \sys{} improves open-vocabulary 3D
semantic segmentation from 52.0 to 58.2 mIoU over the strongest point-only baseline, answers queries in milliseconds
over a relation several times smaller than the raw scan, and \TBD{cracking / workload result}.
\end{abstract}

\section{Introduction}
Point clouds are a first-class data type in robotics, construction, facility management and autonomous driving, but
data systems treat them as opaque arrays. The perception community has recently produced two complementary families of
models: \emph{promptable segmenters} that return the mask of an object from a few clicks~\cite{snap,agile3d,pointsam},
and \emph{3D vision--language models} that attach language-aligned features to every point~\cite{openscene}. Neither is
a query system: promptable segmenters only answer ``what is under this click'', and dense language features only answer
``how similar is each point to this phrase''. Users, however, ask for \emph{objects} that satisfy \emph{predicates}.

\sys{} closes this gap with three contributions.
\begin{itemize}
\item \textbf{A probabilistic, open-vocabulary data model for 3D scenes} (Sec.~\ref{sec:model}). A scene is
materialised once into a relation of object tuples. Semantics are not a fixed label column but a calibrated
conditional probability $P(t\mid o)$ computed for any text $t$ at query time from two embeddings per object.
\item \textbf{\ovql{}, a query language with probabilistic semantics} (Sec.~\ref{sec:lang}) combining text, pointing and
spatial predicates with tuple-independence composition, so answers are ranked by probability.
\item \textbf{Query-driven refinement (semantic cracking)} (Sec.~\ref{sec:proc}). The ingest-time object set is
necessarily incomplete (flat objects on walls, rare concepts). The planner detects predicates no stored object
explains and materialises the missing object by turning the query text into clicks for the promptable segmenter; the
new tuple is persisted with provenance, so the relation adapts to the workload.
\end{itemize}
We also report negative findings that inform the design: fine-tuning the promptable segmenter on a single dataset
loses the cross-dataset generalisation of its multi-dataset training, and multi-mask outputs give little benefit
when selected by a learned quality head (Sec.~\ref{sec:exp}).

\section{Background and Related Work}
\textbf{Promptable 3D segmentation.} Interactive point-cloud segmenters~\cite{interobject3d,agile3d,interactive4d}
return a mask from user clicks; SNAP~\cite{snap} extends this to any domain (indoor, outdoor, aerial) with a single
model and a language-aligned output token. We use SNAP frozen as the geometric expert.
\textbf{Open-vocabulary 3D understanding.} OpenScene~\cite{openscene} distils 2D vision--language features into a 3D
network that runs on points alone; OpenMask3D~\cite{openmask3d} and related work pool image features into class-agnostic
3D masks. These systems produce features or per-scene predictions, not a queryable representation.
\textbf{Data management.} Our probabilistic semantics follow tuple-independent probabilistic databases~\cite{suciu};
query-driven refinement is inspired by database cracking~\cite{cracking}, which reorganises physical data as a
side-effect of queries; vector search~\cite{vectordb} and spatial query processing provide the execution primitives.
\TBD{expand: multimedia/video query systems, LLM-based data systems, spatial DB references}.

\section{Data Model}\label{sec:model}
A scene is voxelised at 2\,cm into $N$ voxels and materialised into
\[
\textsf{Objects}(oid,\; V\subseteq\{1..N\},\; bbox,\; c,\; conf,\; e\in\mathbb{R}^{512},\; p\in\mathbb{R}^{D},\; prov).
\]
$V$ is the object's voxel set, $bbox$ and centroid $c$ support spatial predicates, $conf$ is the segmenter's predicted
mask IoU, $e$ is the segmenter's language token and $p=\mathrm{norm}(\sum_{i\in V} w_i d_i)$ pools the dense 3D
vision--language features $d_i$ inside the mask (mask-anchored pooling). $prov\in\{\texttt{auto},\texttt{click},
\texttt{query:}t\}$ records how the tuple was created. For a phrase $t$ and a background vocabulary $B$,
\[
\log P(t\mid o) \propto \alpha \log \sigma_\tau\!\big(p^\top T_d(t,B)\big) + (1-\alpha)\log P_{e}(t\mid e),
\]
a product of experts in which $P_e$ reads the segmenter token \emph{through its training vocabulary} (vocabulary
relay): the posterior over training anchors is combined with anchor--query text similarity. Ingest runs the frozen
backbone and dense network once, generates automatic prompts, decodes and deduplicates masks.

\section{The \ovql{} Language}\label{sec:lang}
\begin{lstlisting}
SELECT [WHERE expr] [LIMIT k] [THRESHOLD tau]
expr   := term (OR term)*      term := factor (AND factor)*
factor := NOT factor | (expr) | pred
pred   := IS('text') | CONTAINS(x,y,z) | NEAR(expr, d)
        | ABOVE(expr) | BELOW(expr) | ON(expr[, d])
\end{lstlisting}
Every predicate maps an object to a probability. $\textsf{IS}(t)(o)=P(t\mid o)$;
$\textsf{CONTAINS}(x)(o)=conf_o\cdot[x\in V_o]$; for a spatial relation $R$,
$R(E)(o)=\max_{o'\neq o} R(o,o')\,E(o')$ over bounding-box relations; $\textsf{AND}$ is the product, $\textsf{OR}$ the
noisy-or and $\textsf{NOT}$ the complement. Example:
\texttt{SELECT WHERE IS('chair') AND NEAR(IS('table'), 0.5) LIMIT 5}.

\section{Query Processing}\label{sec:proc}
\textbf{Two-tier execution.} The backbone and dense features are computed once per scene; every query touches only
the small decoder (for refinement) and vector arithmetic over $|\textsf{Objects}|$ rows.
\textbf{Phase 1: coverage and semantic cracking.} For each $\textsf{IS}(t)$ in the plan, if $\max_o P(t\mid o)<\theta$
no stored object explains $t$. Promptable segmenters are trained to follow simulated users who click in the error region
of the current mask; we replace the user with the text relevance map $r_t$: the first click is the densest
high-relevance voxel and each further click is the most relevant voxel inside the target region not yet covered. If
geometry cannot separate the object (low mask--relevance agreement, typical of flat objects on walls), the relevance
region itself is used. The tuple is appended with $prov=\texttt{query:}t$ and the decision is recorded (including
negative outcomes), so repeated or related queries do not pay again.
\textbf{Phase 2: probabilistic evaluation} over the now fixed relation, ranking by probability.
\textbf{Verification.} Optionally, a name is accepted only if \textsf{IS}(name) re-segments the same object.

\section{Implementation}
\sys{} uses SNAP-C~\cite{snap} as the frozen promptable segmenter (XYZ-only input, as in its training), the 3D
distilled OpenScene network~\cite{openscene} for dense features, CLIP~\cite{clip} text encoders, and an optional small
vision--language model (Qwen2.5-VL-3B) that chooses among 3D-proposed names from object renders for interactive naming.
\TBD{code size, interactive viewer, artifact link}.

\section{Evaluation}\label{sec:exp}
\textbf{Setup.} ScanNet~\cite{scannet} val (312 scenes, 20 classes) and ScanNet200 splits; S3DIS~\cite{s3dis} Area~5
(zero-shot); Matterport3D~\cite{matterport} 160-class test (406 regions). Evaluation code and protocols are released.

\begin{table}[t]\centering\caption{Open-vocabulary 3D semantic segmentation (mIoU).}\label{tab:semseg}
\begin{tabular}{lccc}\toprule
method & ScanNet & MP-160 & MP-160 unseen \\\midrule
segmenter text head & 41.5 & 6.9 & 1.9 \\
OpenScene 3D (reproduced) & 52.0 & 9.5 & 5.5 \\
\sys{} & \textbf{58.2} & \textbf{11.6} & \textbf{5.7} \\\bottomrule
\end{tabular}\end{table}

\textbf{Open-vocabulary semantics (Table~\ref{tab:semseg}).} Mask-anchored pooling and expert fusion improve over the
point-wise dense baseline by 6.2 mIoU on ScanNet and 2.1 on Matterport-160; on the 63 Matterport classes absent from the
segmenter's training vocabulary \sys{} is on par with the dense baseline (5.7 vs 5.5).

\begin{table}[t]\centering\caption{\ovql{} workloads: instance AP of concept (\textsf{IS}) and relational
(\textsf{IS}$\wedge$\textsf{NEAR}) queries.}\label{tab:queries}
\begin{tabular}{lcccc}\toprule
 & \multicolumn{2}{c}{S3DIS A5} & \multicolumn{2}{c}{ScanNet200 val} \\
engine & IS & NEAR & IS & NEAR \\\midrule
segmenter token only & 9.8 & 12.2 & 8.2 & 29.2 \\
\sys{} & 11.6 & 14.5 & 9.7 & 29.1 \\
\sys{} + cracking & 11.6 & 13.8 & 8.4 & 27.0 \\
\sys{} + cracking (calibrated) & \TBD{} & \TBD{} & \TBD{} & \TBD{} \\\bottomrule
\end{tabular}\end{table}

\textbf{Query workloads (Table~\ref{tab:queries}).} Answer sets are derived from instance annotations; AP is averaged
over IoU 0.50:0.95. On ScanNet200 the segmenter was trained on all 200 classes, so its token is supervised in-domain;
\sys{} improves head classes by 4.7 AP but trails on tail classes (2.5 vs 4.4). Uncalibrated cracking adds false
positives; \TBD{calibrated cracking result}.

\textbf{Interactive segmentation.} With our protocol the released segmenter reproduces its published S3DIS full-room
result (IoU@1/3/5: 53.3/69.1/74.0 vs 53.6/71.1/77.6). Decoder fine-tuning on ScanNet alone degrades zero-shot
multi-click IoU (IoU@5 71.0 vs 72.8 on a held-out subset), and a frozen segmenter with additive single-click tokens has
an oracle headroom of only +2.0 IoU on ScanNet and +3.0 on S3DIS that a learned quality head does not capture; \sys{}
therefore keeps the segmenter frozen.

\textbf{Systems.} \TBD{ingest time vs scene size; query latency on the stored relation ($\approx$2--4\,ms in pilot rooms)
vs re-processing per query ($\approx$6--13\,s); storage; cracking amortisation over a workload with synonyms and absent
terms}.

\textbf{Interactive naming.} \TBD{click-to-name accuracy on Matterport-160 (re-run with XYZ-only input)}.

\section{Conclusion}
\sys{} treats a 3D scan as a relation of objects with open-vocabulary probabilistic semantics, makes it queryable with
a small declarative language, and lets the workload refine the relation. \TBD{limitations and future work: multi-scene
indexes, outdoor domains, learned calibration}.

\section*{Acknowledgment}
\TBD{AI-generated content disclosure as required by ICDE 2027: system(s) used, sections affected, extent.}

\bibliographystyle{IEEEtran}
\begin{thebibliography}{99}
\bibitem{snap} A.~Gupta \emph{et al.}, ``SNAP: Towards segmenting anything in any point cloud,'' in \emph{3DV}, 2026, arXiv:2510.11565.
\bibitem{agile3d} Y.~Yue \emph{et al.}, ``AGILE3D: Attention guided interactive multi-object 3D segmentation,'' in \emph{ICLR}, 2024.
\bibitem{pointsam} Y.~Zhou \emph{et al.}, ``Point-SAM: Promptable 3D segmentation model for point clouds,'' 2024.
\bibitem{interobject3d} T.~Kontogianni \emph{et al.}, ``Interactive object segmentation in 3D point clouds,'' in \emph{ICRA}, 2023.
\bibitem{interactive4d} I.~Fradlin \emph{et al.}, ``Interactive4D: Interactive 4D LiDAR segmentation,'' 2024.
\bibitem{openscene} S.~Peng \emph{et al.}, ``OpenScene: 3D scene understanding with open vocabularies,'' in \emph{CVPR}, 2023.
\bibitem{openmask3d} A.~Takmaz \emph{et al.}, ``OpenMask3D: Open-vocabulary 3D instance segmentation,'' in \emph{NeurIPS}, 2023.
\bibitem{clip} A.~Radford \emph{et al.}, ``Learning transferable visual models from natural language supervision,'' in \emph{ICML}, 2021.
\bibitem{suciu} D.~Suciu, D.~Olteanu, C.~R\'e, and C.~Koch, \emph{Probabilistic Databases}. Morgan \& Claypool, 2011.
\bibitem{cracking} S.~Idreos, M.~L.~Kersten, and S.~Manegold, ``Database cracking,'' in \emph{CIDR}, 2007.
\bibitem{vectordb} \TBD{vector database reference}.
\bibitem{scannet} A.~Dai \emph{et al.}, ``ScanNet: Richly-annotated 3D reconstructions of indoor scenes,'' in \emph{CVPR}, 2017.
\bibitem{s3dis} I.~Armeni \emph{et al.}, ``3D semantic parsing of large-scale indoor spaces,'' in \emph{CVPR}, 2016.
\bibitem{matterport} A.~Chang \emph{et al.}, ``Matterport3D: Learning from RGB-D data in indoor environments,'' in \emph{3DV}, 2017.
\end{thebibliography}
\end{document}
